\chapter*{ECSA self-assessment}

\renewcommand{\arraystretch}{1.5}
\begin{center}
\begin{longtable}{|m{\textwidth}|}
\hline

\textbf{GA 1. Problem solving} \\
\hline    
The project required problem solving to overcome multiple challenges in producing accurate and repeatable hydrogel prints. These included achieving controlled and stable extrusion of gels with complex non-Newtonian rheological properties, ensuring the system was adaptable to a range of bioinks, and providing a method to determine optimal printing parameters for different inks. Problem-solving ability is demonstrated in the mechanical, electronic, firmware, and experimental design aspects of the project. \\
\hline    

\textbf{GA 2. Application of scientific and engineering knowledge} \\
\hline    
The project required the integration of knowledge across mechanics, electronics, biology, and experimental design to develop a functional hydrogel printer. Mechanical analysis was used to calculate buckling loads in the syringe and to size actuators by determining the forces and speeds required for each stepper motor. Fluid mechanics principles, including Poiseuille’s law, were applied to estimate extrusion pressures and inform the syringe design. Circuit analysis and embedded systems knowledge were applied in designing and programming the custom PCB for control and sensing. Biological understanding of tissue engineering provided the performance requirements that shaped the system’s capabilities, while experimental design ensured validation of the printer's capabilities. \\
\hline    

\textbf{GA 3. Engineering Design} \\
\hline    
Engineering design was a major aspect of the project and was carried out across three primary phases: mechanical, embedded system, and experimental design. The mechanical design focused on developing the syringe-based extruder and gantry; the embedded system design centred on creating a custom PCB to control motion and sensing; and the experimental design established a methodology for determining the optimal printing parameters of a chosen bioink. Together these phases demonstrate a structured, multidisciplinary approach to achieving a functional bioprinter. \\
\hline    

\textbf{GA 5. Engineering methods, skills and tools, including Information Technology} \\
\hline    
The project required a wide range of engineering methods, skills and tools to design, build, and evaluate the 3D printer. Inventor Professional was used to create a 3D model of the mechanical bioprinter design. Analytical methods were applied to calculate mechanical loads, extrusion pressures, and motor requirements. Embedded hardware was designed using KiCad and the firmware was created through STM32CubeIDE. Experimental methods, supported by data processing in MATLAB, were applied to test the system and optimise printing parameters. \\
\hline    
    
\textbf{GA 6. Professional and technical communication} \\
\hline    
The project outcomes were communicated through an oral presentation and a detailed written report, supported by CAD models, schematics, and results. These demonstrated the ability to communicate information regarding technical work to engineering audiences. \\
\hline    

\textbf{GA 8. Individual, Team and Multidisciplinary Working} \\
\hline    
The project demonstrated independent execution while drawing on knowledge from multiple disciplines, including mechanics of materials, fluid mechanics, electronics, embedded software development, and biology. Collaboration with Stellenbosch University’s workshop artisans and the Physiology Department ensured the system was manufacturable and aligned with clear design objectives. \\
\hline    

\textbf{GA 9. Independent Learning Ability} \\
\hline    
Independent learning was demonstrated in several aspects of the project such as the literature review, which provided foundational knowledge in tissue engineering and hydrogel material science, as well as exposure to new approaches in mechanical design. New skills were gained in PCB design using KiCad, which enabled the development of the bioprinter's circuit board for sensing and control. \\
\hline    

\end{longtable}
\end{center}